% Transcription — fonds Alexandre Grothendieck, shelfmark 56, batch 1 (pages 1-20).
% Established from the facsimile archives/batches/56/batch-01.pdf (a mirror of
% https://grothendieck.umontpellier.fr/56.pdf), per the skill
% transcribe-grothendieck. The folder is thirty-eight pages; this is the first
% of two batches (pages 21-38 are batch 2).
%
% Pages skipped: 9 and 18, blank cream sheets with nothing on them. Pages 1 and
% 8 are tan covers inscribed in his hand; they take no \page marker, and their
% words become the two section headings, a \note saying where each comes from.
% Page 1 reads « Situation des "vanishing cycles" : [one illegible word] Cas des
% quadriques », with a pencilled « 3 p » that is not obviously his and is not
% reproduced; page 8 reads « Calculs de cohomologie des quadriques, des cônes
% projetants », with a pencilled « 12 p », likewise left out.
%
% The manuscript notes (pages 2, 4, 6, 10, 12, 14, 16, 19) are written on the
% backs of leaves of a typescript of his own: pages 3, 5, 7, 11, 13, 15, 17 and
% 20 are leaves of « XII Tores maximaux, groupe de Weyl, sous-groupes de Cartan,
% centre réductif des schémas en groupes lisses et affines, par A.
% Grothendieck » — to all appearances the typescript of what was published as
% Exposé XII of SGA 3 — corrected by him in ink. The typescript is unrelated to
% the notes on its versos. Following folders 48, 54 and 55, and the brief for
% this folder (the typed text is his, but published), the typed text is not
% re-set: each typescript page gets one French \note summarising what it states,
% and his interventions on it (insertions, substitutions, deletions, marginal
% notes, the inked renumbering of the paragraphs as 1.x and 2.x) are given in
% full. These eight pages are therefore « summarised », though the text is his.
% Several of them are crossed by long pencil strokes over the whole page (3, 5)
% or by ink crosses over a cancelled passage (11, 15); recorded once per page.
%
% What the batch is about. Two runs of notes on the local monodromy of a
% degeneration to a quadric with one singular point — the « vanishing cycles »
% of the inventory title — in ℓ-adic (étale) cohomology with coefficients μ.
% Pages 2-6 (under the first cover): X → S projective flat over a strictly
% local discrete valuation ring, special fibre X_0 a quadric with a singular
% point, geometric generic fibre \bar X_1 a smooth quadric, fibre dimension
% n-1; the question « H^i(X_0) → H^i(\bar X_1) ? » and the exact sequence
% ... → H^i(X_0) → H^i(\bar X_1) → Φ^{i+1} → H^{i+1}(X_0) → ... defining the
% Φ^i as kernels and cokernels of specialisation. Page 4 writes out the
% cohomology rings of both fibres in terms of the hyperplane class ξ and the
% extra classes λ'_ν (in X_0 for n even) or λ_ν (in \bar X_1 for n odd), and
% concludes in the two parities: Φ^i = 0 for i ≠ n, Φ^n ≃ μ(-ν) or μ(...)
% (« isom. déterminé au signe près »), R^i f_* constant for i ≠ n resp. n-1,
% and π acting trivially on H^*(\bar X_1) for n even but not on Φ^n for n odd.
% Page 6, « Calcul affine », replaces X by U = X - Z (Z at infinity) and gets
% directly H^i(\bar U_1) ≃ Φ^{i+1}, Φ^i = 0 for i ≠ n, Φ^n = H^{n-1}(\bar U_1)
% generated by the « vanishing cycle ».
% Pages 10-19 (under the second cover): cohomology of a projective cone.
% Page 10: the diagram of C ⊂ \hat C ⊃ X (cone, projective cone, section at
% infinity) and its primed copy, E = C - a, \hat E = \hat C - a; H^i(C) = 0 for
% i ≠ 0, the short exact sequences, H^i(\hat C) ≃ H^{i-2}(X)(-1) for i ≠ 0 via
% the Gysin map of j : X → \hat C, j_*(a)j_*(b) = j_*(abξ), and the Gysin
% sequence of the line bundle \hat E → X. Page 12: for X smooth connected of
% dimension n, H^i(E) ≃ P^i(X) mod torsion for i ≤ n and ≃ P^{i-2}(X)(-1) for
% i ≥ n+1 « en vertu du th. de Lefschetz »; the local cohomology at the vertex,
% H^0_a(C) = H^1_a(C) = 0, H^i_a(C) ≃ H^{i-1}(E) for i ≥ 2; and for X a
% quadric Q^n, \hat C = Q'^{n+1}, the rings of Q'^{n+1} in the two parities
% with ξλ'_{ν+1} = 0, λ'^2_{ν+1} = 0. Page 14: the resulting table of H^i(E)
% for the cone over a quadric (Z/2Z appears in degree 2ν-2 for ℓ = 2 when
% n+1 = 2ν). Page 16: cohomology with compact supports of C, E, \hat E, ending
% on the boxed question « H^*_a(C) ≃ H^*_c(C) ?? ». Page 19, « Cohomologie des
% quadriques »: Q^n smooth over k separably closed, Q^{n-1} its section at
% infinity, U^n = AQ^n = Q^n - Q^{n-1}; H^i_c(U^n, A) and H^i(U^n, A) with the
% twists ν_n, ν'_n according to the parity of n; the map φ from compact-support
% to ordinary cohomology, zero except for n even, i = n; the Gysin sequence
% H^{i-2}(Q^{n-1})(-1) → H^i(Q^n) → H^i(U^n) → ..., surjective for i ≠ 0, n,
% injective for i ≠ 1, n+1; it breaks off in the case n even.
%
% The hand of the manuscript pages is fast drafting, and on pages 2 and 16 in
% particular most of the connective prose does not carry; following folder 29
% batches 8-9 and the passes since, it is left \ill{} and the formulas are
% given in full. Pages 4, 10, 12, 14 and 19 read much better.
%
% Notation fixed once for the batch. The fibres are X_0 and \overline{X}_1;
% the letter read as Φ (a barred circle) for the vanishing-cycle groups; μ for
% the sheaf of roots of unity he writes with a doubled stem, μ(1), μ(-ν);
% ξ the hyperplane class; a the vertex of the cone, read as the index of the
% local cohomology H^i_a (the subscript could also be x); \widehat{C} the
% projective cone. Degrees and twists overwritten on page 14 are given in their
% final form, with a note.
%
% Readings to check first: the word after « Situation des "vanishing cycles" »
% on page 1 (\ill); « morphisme » (page 2, first line); the whole N.B. of
% page 2; « Calcul affine » (page 6); the second word of the heading of page 10;
% the superscripts of P on page 12 and of H on page 14 (overwritten); the
% twists ν_n, ν'_n on page 19 and the margin note there, read turned
% (« on ne sait déterminer qu'au signe près »).
%
% Pass: Opus 5.5 (claude-opus-5-5), 2026-09-24 — first pass, unchecked against the pages by a human.

\documentclass[11pt,a4paper]{article}
% Le préambule commun à toutes les transcriptions du fonds Grothendieck.
%
% Il tient en une idée : **l'incertitude se compose, elle ne se résout pas.**
% Une transcription qui lisse un mot douteux a détruit la seule chose pour
% laquelle on la faisait — un lecteur doit pouvoir savoir, à chaque endroit,
% ce qui était sur le feuillet et ce qui a été ajouté. Les six macros qui
% suivent sont donc l'appareil critique, et non de la décoration.
%
% Les mêmes commandes sont reconnues par `scripts/render.mjs`, qui produit la
% vue de lecture du site. Une macro ajoutée ici doit l'être aussi là-bas,
% faute de quoi le rendu échouera bruyamment — ce qui est voulu : un rendu
% silencieusement faux est pire qu'un rendu absent.

\ProvidesPackage{grothendieck}[2026/08/08 Transcriptions du fonds Grothendieck]

\usepackage{amsmath,amssymb,amsthm}
\usepackage{mathrsfs}
\usepackage{stmaryrd}
% Les diagrammes commutatifs sont partout dans ce fonds : c'est la langue
% même de la géométrie algébrique de Grothendieck, et une transcription qui
% les rendrait en prose ne transcrirait rien.
\usepackage{tikz-cd}
% Français par défaut — c'est la langue des feuillets — mais l'anglais est
% chargé aussi : la traduction et le résumé sont des documents anglais, et la
% typographie française (espaces avant les deux-points) y serait une faute.
% Ces documents basculent par \selectlanguage{english} en tête de corps.
\usepackage[english,french]{babel}
\usepackage{fontspec}
\usepackage{geometry}
\geometry{a4paper,margin=25mm,marginparwidth=28mm}
\usepackage{marginnote}
\usepackage{xcolor}
% ulem, not soul. `soul`'s \st and \ul are built on a character-by-character
% scan that XeLaTeX's font machinery breaks: the document fails with an
% undefined control sequence rather than degrading. ulem does the same two jobs
% and is Unicode-engine safe. `normalem` stops it hijacking \emph, which the
% transcriptions use for Grothendieck's own emphasis.
\usepackage[normalem]{ulem}
\usepackage{hyperref}
\hypersetup{colorlinks=true,linkcolor=black,urlcolor=black,pdfborder={0 0 0}}

% --- Métadonnées du lot -----------------------------------------------------
% Reprises telles quelles par le script de rendu, qui les affiche en tête de
% la vue de lecture. Elles fixent ce que le fichier prétend couvrir : sans
% elles, une transcription flotte sans point d'attache dans un fonds de seize
% mille feuillets.

\newcommand{\folder}[1]{\gdef\grothFolder{#1}}
\newcommand{\batch}[1]{\gdef\grothBatch{#1}}
\newcommand{\leaves}[2]{\gdef\grothFirst{#1}\gdef\grothLast{#2}}
\newcommand{\foldertitle}[1]{\gdef\grothFoldertitle{#1}}
\gdef\grothFolder{?}\gdef\grothBatch{?}\gdef\grothFirst{?}\gdef\grothLast{?}\gdef\grothFoldertitle{}

% --- Le résumé --------------------------------------------------------------
%
% Il ouvre la lecture modernisée, sous le titre « Résumé », et il est composé
% en petit. Ce n'est pas un ornement : c'est le seul passage du document qui
% ne soit pas des mathématiques, et un lecteur doit pouvoir le reconnaître
% avant de l'avoir lu — puis le sauter s'il connaît déjà le sujet, ou s'y
% arrêter s'il ne le connaît pas. Le filet qui le suit marque la reprise du
% corps.

\newenvironment{resume}
  {\small\setlength{\parskip}{0.45em}}
  {\par\medskip\hrule\medskip}

% --- Mots-clés ---------------------------------------------------------------
%
% En anglais, en clôture du résumé de la lecture modernisée : le vocabulaire
% moderne sous lequel chercher ce que les feuillets construisent. Le manifeste
% du site (`npm run manifest`) les extrait pour étiqueter la cote entière —
% cette ligne est la source unique des tags, il n'y a pas de fichier à côté.
\newcommand{\keywords}[1]{\par\smallskip\noindent
  {\footnotesize\textsc{Keywords} --- \textit{#1}}\par}

% --- L'appareil critique ----------------------------------------------------

% Le numéro de feuillet, en marge. C'est l'ancre qui permet de régler un
% désaccord sur une formule en regardant *un* feuillet plutôt qu'une chemise
% de six cents. Le site s'en sert aussi pour tourner les pages du fac-similé
% quand on fait défiler la transcription.
\newcommand{\leaf}[1]{%
  \marginnote{\footnotesize\color{gray!70}\textsf{#1}}%
  \label{leaf:#1}\ignorespaces}

% Illisible : on ne devine pas. Un mot inventé qui se lit comme les autres est
% le pire résultat possible.
\newcommand{\ill}{\textcolor{red!60!black}{[\,\ldots\,]}}

% Lecture douteuse : le mot est donné, et signalé comme incertain.
\newcommand{\uncertain}[1]{\uline{#1}}

% Ajout de l'éditeur — une lettre manquante, un mot sous-entendu.
\newcommand{\add}[1]{\textcolor{blue!50!black}{[#1]}}

% Biffé par Grothendieck lui-même. Ce qu'il a rayé fait partie du document :
% une rature dit souvent où la pensée a changé de direction.
\newcommand{\struck}[1]{\sout{#1}}

% Note du transcripteur, sur l'état matériel du feuillet.
\newcommand{\note}[1]{\par\smallskip{\small\color{gray!60!black}\textit{[#1]}}\par\smallskip}

% Note marginale de Grothendieck, distincte du corps du texte.
\newcommand{\marginal}[1]{\par\smallskip\noindent\hspace{1em}%
  {\small\color{gray!60!black}#1}\par\smallskip}

% --- Titre ------------------------------------------------------------------

\AtBeginDocument{%
  \begin{center}
    {\large\bfseries Fonds Alexandre Grothendieck --- cote n\textsuperscript{o}~\grothFolder}\\[2pt]
    {\itshape\grothFoldertitle}\\[4pt]
    {\small feuillets \grothFirst--\grothLast\ (lot \grothBatch)}\\[2pt]
    {\footnotesize\color{gray!60!black}Transcription établie d'après le fac-similé numérisé
     par l'Université de Montpellier.\\
     Les lectures incertaines sont soulignées, les illisibles notées [\,\ldots\,],
     les ajouts entre crochets.}
  \end{center}
  \vspace{1em}}

\endinput


\folder{56}
\batch{1}
\pages{1}{20}
\dating{1962}
\watermark{Édition de démonstration}
\foldertitle{Lefschetz-Hodge - Vanishing cycles : notes manuscrites (s.d.), tiré à part (1962)}

\begin{document}

\section*{Situation des « vanishing cycles » : \ill{} Cas des quadriques}

\note{titre de sa main, en haut à droite du feuillet de garde (page 1), qui ne
porte rien d'autre ; c'est lui qui nomme la suite, et le feuillet ne reçoit pas
de numéro de page}

\page{2}

\marginal{Vanishing cycles \uncertain{: cas} des quadriques}\note{inscrit en
biais dans le coin supérieur gauche, séparé du reste de la page par un trait
oblique}

\begin{tikzcd}
X \arrow[d, "f"] \\
S
\end{tikzcd}

$f$ \uncertain{morphisme} projectif plat

$S$ spectre d'un anneau de val. discrète strictement local

$X_0$ fibre spéciale quadrique à point singulier

$\overline{X}_1$ fibre générique géométrique, quadrique non singulière

dimension \uncertain{commune} des fibres $n-1$

\struck{\ill{}}

$X \subset \mathbf{P}^{n}_{S}$ défini par une équation quadratique …

$H^i(X_0) \to H^i(\overline{X}_1)$ ? (\uncertain{Cf. la suite} exacte générale
\[
  \cdots \to H^i(X_0) \to H^i(\overline{X}_1) \to \Phi^{i+1} \to H^{i+1}(X_0)
  \to H^{i+1}(\overline{X}_1) \to \cdots
\]

\uncertain{On sait par l'}\ill{} \ill{} de \struck{$H^i(X_0)$ et}\note{ligne
portée en interligne au-dessus de la suivante}

1°) \emph{$n$ pair}, \struck{\ill{} $n$ impair}, $X_0$ et $\overline{X}_1$ ont
\ill{} \ill{} en dim \uncertain{impaire}. Donc \ill{} \ill{}
\[
  0 \to \Phi^{2\nu} \to H^{2\nu}(X_0) \to H^{2\nu}(\overline{X}_1) \to
  \Phi^{2\nu+1} \to 0
\]
i.e. les $\Phi^i$ \ill{} \ill{} \ill{} les noyaux et conoyaux du \ill{} de
\emph{spécialisation}.

N.B. \ill{} \ill{} \ill{} \ill{} $X$ \ill{} $X_0$, \ill{} \ill{} \ill{}
$\xi \in H^2(X, \mu(1))$, \uncertain{provenant} de la \uncertain{section}
\ill{} \ill{} \ill{} \ill{} \ill{} \ill{} \uncertain{restriction} \ill{} \ill{}
\uncertain{fibre} \ill{} …\note{les quatre dernières lignes de la page sont
d'une écriture très rapide ; seuls $X$, $X_0$ et la classe $\xi$ se lisent sûrement}

\page{3}

\note{feuillet de tapuscrit, repris à la main, de l'exposé « XII Tores
maximaux, groupe de Weyl, sous-groupes de Cartan, centre réductif des schémas
en groupes lisses et affines » (voir page 7), sans rapport avec les notes
manuscrites au verso desquelles il se trouve ; il est gardé parce qu'il porte
sa main. Le texte tapé est le sien et publié : on en donne la teneur et ses
interventions, sans le recomposer. La page est barrée de trois longs traits de
crayon. Elle donne les énoncés a) à d) d'un théorème sur les fonctions
$\rho_r$ (rang réductif), $\rho_n$ (rang nilpotent) et $\rho_u = \rho_n - \rho_r$
sur $S$ : a) $\rho_r$ est semi-continue inférieurement, $\rho_n$ et $\rho_u$
supérieurement ; b) équivalence de (i) $\rho_r$ localement constante, (ii)
existence locale, pour la topologie étale, d'un tore maximal dans $G$, (ii bis)
la même chose pour la topologie fidèlement plate quasi-compacte ; c) deux tores
maximaux sont localement conjugués pour la topologie étale ; d) sous b),
$\rho_n$ (donc $\rho_u$) est aussi localement constante ; puis le début de la
démonstration de a), par réduction au cas d'un changement de base fidèlement
plat quasi-compact (SGA VIII 4.3)}

\begin{itemize}
  \item les lettres $\rho$ (et $\rho_r$, $\rho_n$, $\rho_u$, $\rho'$) sont
  partout portées à la main dans les blancs laissés par la machine ;
  \item « localement conjugués » est entouré d'un trait ;
  \item « la topologie \add{de Zariski} de $S$ étant en effet quotient de celle
  de $S'$ » : « de Zariski » inséré à la main au-dessus de la ligne, après
  « \struck{(SGA VIII)} », biffé.
\end{itemize}

\page{4}

1°) \emph{$n$ pair $= 2\nu$}, \emph{$n-1$ impair} \quad $n-1 = 2\nu - 1 = 2(\nu-1)+1$

\[
  \left.
  \begin{array}{ll}
    H^{*}(X_0) : & 1, \xi_0, \ldots, \xi_0^{\nu-1}, [\xi_0^{\nu}, \lambda'_{\nu}],
      \tfrac{1}{2}\xi_0^{\nu+1}, \ldots, \tfrac{1}{2}\xi_0^{2\nu-1} \\[4pt]
    H^{*}(\overline{X}_1) : & 1, \xi_1, \ldots, \xi_1^{\nu-1}, \tfrac{1}{2}\xi_1^{\nu},
      \tfrac{1}{2}\xi_1^{\nu+1}, \ldots, \tfrac{1}{2}\xi_1^{2\nu-1}
  \end{array}
  \right]
  \quad
  \begin{array}{l}
    \xi_0 \mapsto \xi_1 \\
    \lambda'_{\nu} \mapsto 0
  \end{array}
\]

\emph{Donc} $H^i(X_0) \to H^i(\overline{X}_1)$ \emph{bijectif} \struck{excepté}
si $i \neq 2\nu = n$

$H^{2\nu}(X_0) \to H^{2\nu}(\overline{X}_1)$ \emph{surjectif}, noyau engendré par
$\lambda'_{\nu}$

\emph{Donc}
\begin{itemize}
  \item[a)] $\Phi^i = 0$ \struck{\ill{}} si $i \neq (2\nu = n)$ ;
  $\Phi^{\uncertain{2\nu}} \simeq \mu(-\nu)$ (isom. déterminé au signe près …)
  \note{l'exposant de $\Phi$ et le $-\nu$ sont repassés et empâtés}
  \item[b)] $R^i f_*$ constant si $i \neq n$ ; $R^n f_*$ = un \ill{}-faisceau
  concentré en $s$, isom. à $\mu(\ldots)$.
  \marginal{\ill{} \ill{} \ill{} en \uncertain{dimension} \uncertain{constante}}
  \item[c)] $\pi$ opère trivialement sur $H^{*}(\overline{X}_1)$ et sur
  $\Phi^{2\nu}$\note{c) est écrit dans la marge gauche, en regard de b)}
\end{itemize}

2°) \emph{$n$ impair} $= 2\nu+1$, \quad $n - 1 = 2\nu$

\[
  \left.
  \begin{array}{ll}
    H^{*}(X_0) : & 1, \xi_0, \ldots, \xi_0^{\nu-1}, \xi_0^{\nu},
      \tfrac{1}{2}\xi_0^{\nu+1}, \ldots, \tfrac{1}{2}\xi_0^{2\nu} \\[4pt]
    H^{*}(\overline{X}_1) : & 1, \xi_1, \ldots, \xi_1^{\nu-1}, [\xi_1^{\nu}, \lambda_{\nu}],
      \tfrac{1}{2}\xi_1^{\nu+1}, \ldots, \tfrac{1}{2}\xi_1^{2\nu}
  \end{array}
  \right]
  \quad \xi_0 \mapsto \xi_1
\]

Donc $H^i(X_0) \to H^i(\overline{X}_1)$ bijectif si $i \neq 2\nu = n-1$

$H^{2\nu}(X_0) \to H^{2\nu}(\overline{X}_1)$ injectif, conoyau engendré par
\uncertain{$\lambda_{\nu}$}\note{le dernier symbole est surchargé au point
d'être à peine lisible}

\begin{itemize}
  \item[a)] $\Phi^i = 0$ si $i \neq 2\nu + 1 = n$ ; $\Phi^n \simeq \mu(\ldots)$
  \item[b)] $R^i f_*$ constant si $i \neq n-1$ ; $R^{n-1} f_*$ est « pur » …
  \item[c)] $\pi$ opère trivialement sur $H^i(\overline{X}_1)$ pour
  $i \neq n-1$, mais \uncertain{non} sur $\Phi^n$ …\note{écrit dans la marge
  gauche, au pied de la page}
\end{itemize}

\page{5}

\note{feuillet de tapuscrit du même exposé XII, paginé « 1 » : le début du
paragraphe 1, « Tores maximaux ». Il rappelle qu'on utilisera désormais le
Séminaire Chevalley 56/57 (« BIBLE ») ; définit, pour $G$ algébrique sur $k$
algébriquement clos, les tores maximaux et note qu'ils coïncident avec ceux de
$G_{\text{réd}}$ ; cite le théorème de Borel (conjugaison des tores maximaux,
BIBLE 6, th. 4 c) et définit le \emph{rang réductif} de $G$ ; signale que
l'hypothèse « $G$ affine » est inutile grâce au théorème de Chevalley sur les
extensions d'une variété abélienne par un groupe affine ; commence la
définition du sous-groupe de Cartan. Le haut de la page est couvert d'une
rangée de croix au crayon, et deux longs traits de crayon la barrent de haut en
bas}

\begin{itemize}
  \item le titre : « 1. Tores maximaux \struck{et groupe de Weyl} », le « 1 »
  repassé à l'encre ;
  \item « conjugués par un élément de $G(k)$ \add{$= G_{\text{réd}}(k)$} »,
  inséré au-dessus de la ligne ; en marge, au crayon et entouré,
  « $G_{\text{réd}}(k)$ » ;
  \item « affirmant que tout groupe algébrique connexe lisse sur $k$ est une
  extension d'une variété abélienne par un groupe affine ; \add{cf N° 6}
  \struck{nous n'aurons pas besoin de cette généralisation} » ;
  \item « Dans \add{les Nos 1 à 4}, nous nous bornerons \add{le plus souvent
  aux} \struck{aux groupes} préschémas en groupes affines sur la base. » —
  avec, en marge, « pas à la ligne » ;
  \item « \add{Soit $G$ un groupe algébrique lisse sur $k$,} \struck{$G$ sur
  $k$ étant comme ci-dessus,} et $T$ un tore maximal de $G$, le centralisateur
  \add{$C(T)$} de $T$ dans $G$ sera appelé le » — avec, en marge, « à la
  ligne ».
\end{itemize}

\page{6}

\emph{Calcul \uncertain{affine}} \quad Nous remplaçons $X$ par $U = X - Z$
(\struck{\ill{}} \ill{} à l'infini), on a ainsi
\[
  \cdots \to H^i(U_0) \to H^i(\overline{U}_1) \to \Phi^{i+1} \to H^{i+1}(X_0)
  \to H^{i+1}(\overline{X}_1) \to \cdots
\]
\note{les deux derniers termes portent bien $X$, non $U$}

Or on a maintenant $\underline{H^i(U_0) = 0}$ \quad $i \neq 0$,

d'où
\[
  H^i(\overline{U}_1) \simeq \Phi^{i+1}
\]
d'où directement
\[
  \Phi^i = 0 \ \text{si}\ i \neq n, \qquad
  \Phi^n = H^{n-1}(\overline{U}_1)
\]
engendré par le « vanishing cycle » …

\page{7}

\note{feuillet de tapuscrit paginé « 0 » : la page de titre de l'exposé, « XII
Tores maximaux, groupe de Weyl, sous-groupes de Cartan, centre réductif des
schémas en groupes lisses et affines. par A. Grothendieck ». Le texte tapé
annonce qu'à partir de cet exposé on fera usage des résultats connus sur les
groupes algébriques lisses sur un corps algébriquement clos, notamment de la
théorie de Borel exposée au Séminaire Chevalley 56/57 (« BIBLE », exposés 4 à
7), la théorie des schémas n'y apportant pas de simplification notable ; le
but est d'en déduire des résultats analogues sur une base quelconque, à la
différence de la théorie de structure de Chevalley des groupes semi-simples,
qui se traite ab ovo sur une base quelconque}

\begin{itemize}
  \item « BOREL » et « CHEVALLEY » repassés à la main en capitales, trois fois ;
  \item « Il \struck{\ill{}} semble \add{d'ailleurs} \struck{\ill{}} que la théorie des schémas
  \supplied{n'}apporte aucune simplification notable à la théorie \struck{exposée} de
  BOREL, \struck{\ill{}} telle qu'elle est exposée dans BIBLE. C'est pourquoi il
  n'a pas semblé utile de la reproduire \struck{\uncertain{notre objet étant}} dans le
  présent Séminaire » ;
  \item après « … sur un préschéma de base quelconque), » il ajoute, entouré :
  « (que nous avons suivi dans les Nos 1 à 4) » ;
  \item puis, en bas de page, un paragraphe entier de sa main :
\end{itemize}

\begin{quote}
Dans les exposés oraux, nous nous étions limités aux préschémas en groupes
affines sur $S$, en nous appuyant de façon essentielle sur les théorèmes de
représentabilité de Exp XI N° 4. Dans les présentes notes (cf Nos 6 à 8) nous
montrons rapidement comment on peut éliminer les hypothèses affines
\struck{\ill{} Exp XI \ill{}} par une \struck{autre} méthode plus simple
n'utilisant pas les résultats les plus délicats de Exp. XI.
\end{quote}

\marginal{Pour d'autres généralisations des résultats contenus dans le présent
exposé, voir également exposés \struck{Exp.} XV et XVI. (Il est évident que le
contenu des exposés XI, XII, XV, XVI devrait être complètement
refondu).}\note{écrit de bas en haut dans la marge gauche, lu en tournant la
feuille ; un signe de renvoi le rattache au bas du paragraphe manuscrit}

\section*{Calculs de cohomologie des quadriques, des cônes projetants}

\note{titre de sa main, en haut à droite du feuillet de garde (page 8), qui ne
porte rien d'autre ; il nomme la suite et ne reçoit pas de numéro de page}

\page{10}

\emph{Cohomologie \ill{} d'un cône projectif}

\begin{tikzcd}
X \arrow[d, "h"] & & \\
C' \arrow[r, hook, "i'"'] \arrow[d] \arrow[rr, bend left=35] & \widehat{C}' \arrow[d] & X' \arrow[l, hook', "j^{\prime}"] \arrow[d, "h"] \\
C \arrow[r, hook, "i"'] \arrow[rr, bend right=35] & \widehat{C} & X \arrow[l, hook', "j"]
\end{tikzcd}

\note{le diagramme est redessiné d'après le croquis ; la flèche $h$ de la
colonne de gauche descend de $X$ vers $C'$, telle qu'elle est tracée}

$E = C - a$

$\widehat{E} = \widehat{C} - a$

\[
  \struck{H^{i-1}(X) \to H^i(\widehat{C}) \to H^i(\widehat{C}') \to H^i(X) \to
  H^{i+1}(\widehat{C})}
\]
\struck{$H^i(C) \to$}\note{au-dessus de la suite biffée, « $i \geqslant 0$ »}

\ill{}
\[
  \begin{array}{ccccccc}
    \cdots \to H^{i-1}(X) & \to & H^i(\widehat{C}) & \to & H^i(\widehat{C}') & \to & H^i(X) \to \cdots \\
    \downarrow & & \downarrow & & \downarrow & & \| \\
    H^{i-1}(X) & \to & H^i(C) & \to & H^i(C') & \xrightarrow{\text{bij}} & H^i(X)
  \end{array}
\]

\[
  \boxed{H^i(C) = 0 \quad \text{si}\ i \neq 0}
\]
\[
  0 \to H^i(\widehat{C}) \to H^i(\widehat{C}') \to H^i(X) \to 0
\]
\[
  \boxed{H^i(\widehat{C}) \xleftarrow{\ \sim\ } H^{i-2}(X)(-1) \qquad i \neq 0}
\]
provenant de l'injection \uncertain{de Gysin} $X \xrightarrow{j} \widehat{C}$

$j_*(a)\, j_*(b) = j_*(ab\xi)$ \quad car $j^* j_*(c) = c\xi$, ce qui détermine
la structure multiplicative …

\struck{$E$} $\widehat{C} - a \to$ fibré sur $X$ en fibres des droites
\struck{aff.} vectorielles,

$\simeq \widehat{C}' - a$, donc
\[
  \boxed{H^{*}(X) \xrightarrow{\ \sim\ } H^{*}(\widehat{C} - a)}, \qquad
  \widehat{C} - a = \widehat{E}, \qquad \widehat{C} - a - X = E
\]
\[
  H^{i-2}(X)(-1) \to H^i(\widehat{E}) \to \boxed{H^i(E)} \to H^{i-1}(X)(-1)
  \to H^{i+1}(\widehat{E})
\]
avec $H^i(\widehat{E}) \simeq H^i(X)$, $H^{i+1}(\widehat{E}) \simeq H^{i+1}(X)$ ;
la première flèche est la multiplication par $\xi$ (cup-$\xi$).\note{légende
écrite en biais sous la première flèche}

\page{11}

\note{feuillet de tapuscrit du même exposé : « Remarques 1.11 ». a) le
préschéma $\mathcal{T}$ de 6.10 s'appellera le \emph{préschéma des tores
maximaux} de $G$ ; on verra au N° 5 qu'il est affine sur $S$, ce qu'on peut
vérifier directement quand $G$ est un sous-préschéma en groupes fermé de
$\mathrm{Gl}(n)$, $\mathcal{T}$ s'identifiant alors à un sous-préschéma ouvert
et fermé du préschéma qui représente les sous-groupes de type multiplicatif de
$G$. b) sur une démonstration de 1.7, donc de 1.9, qui n'utilise pas les
résultats délicats de l'exposé XI}

\begin{itemize}
  \item le « s » de « Remarques » et le « 1 » de « 1.11 » sont de sa main ;
  \item « Il \struck{ne serait pas difficile} \add{est possible} de donner une
  démonstration de 1.7. donc de 1.9. n'utilisant pas les résultats délicats de
  Exp XI, mais seulement \add{les résultats faciles} \struck{\ill{}} Exp XI
  3.12. \add{et 6.2.}, en travaillant uniquement avec des groupes de type
  multiplicatif finis sur la base » ; en marge, au crayon et entouré,
  « et 6.2., » ;
  \item « comparer \struck{Exp} N° 7. \struck{Par contre,} \add{La même remarque
  vaut pour} la démonstration de 1.10., » — la suite tapée (« s'appuie de façon
  essentielle sur Exp XI 4.1. Dans un autre exposé, nous donnerons une version
  généralisée de 1.10. par une démonstration différente … les théorèmes de
  représentabilité de Exp XI N° 4 pourraient être éliminés de la théorie
  présentée dans le présent exposé et les suivants ») est encadrée, biffée
  ligne à ligne et barrée de grandes croix, et remplacée par :
\end{itemize}

\begin{quote}
\add{de façon générale,} il est possible (et sans doute avantageux) d'éliminer
totalement les résultats de représentabilité de Exp XI N° 4 de la théorie
présentée dans le présent exposé et les suivants.
\end{quote}
\note{« de façon générale, » est écrit deux fois, dans la marge gauche,
entouré, et au bout de la dernière ligne tapée ; l'encadré manuscrit est
lui-même traversé de quelques traits obliques fins, qui prolongent peut-être
les croix du passage biffé}

\page{12}

Si $X$ est non singulière (\uncertain{connexe de dim $n$}) on trouve :
\[
  \boxed{
  \begin{array}{lll}
    i \leqslant n & H^i(E) \simeq P^i(X) & \text{mod torsion} \\
    i \geqslant n+1 & H^i(E) \simeq P^{\uncertain{i-2}}(X)(-1) & \text{mod. de torsion}
  \end{array}}
\]
\marginal{en vertu du th. de Lefschetz}\note{la lettre $P$ de la seconde ligne est
repassée en gras et son exposant empâté}

\[
  \cdots \to H^{i-1}(C) \to H^{i-1}(E) \to H^i_a(C) \to H^i(C) \to H^i(E) \to
  H^{i+1}_a(C) \to \cdots
\]
avec $H^{i-1}(C) = 0$ si $i - 1 \neq 0$, $H^i(C) = 0$ si $i \neq 0$.

\struck{$i \neq 0, 1$}
\[
  \to H^0_a(C) \to H^0(C) \xrightarrow{\text{bij}} H^0(E) \to H^1_a(C) \to
  H^1(C) \to H^1(E) \to H^2_a(C) \to H^2(C)
\]
avec $H^0_a(C) = 0$, $H^1(C) = 0$, $H^2(C) = 0$ ; la flèche bij si $X$ connexe.

\[
  \boxed{
  \begin{array}{l}
    H^0_a(C) = H^1_a(C) = 0 \\
    H^i_a(C) \simeq H^{i-1}(E) \quad \text{si}\ i \geqslant 2
  \end{array}}
\]

Cas $X$ une quadrique $Q^n$, donc $\widehat{C} = Q'^{\,n+1}$, on trouve
\[
  \underline{H^i(Q'^{\,n+1}) \xleftarrow{\ \sim\ } H^{i-2}(Q^n)(-1)} \quad
  \text{si}\ i \neq 0 \quad \text{donc :}
\]
\[
  \begin{array}{ll}
    \left(1, \xi, \ldots, \xi^{\nu}, \tfrac{1}{2}\xi^{\nu+1}, \ldots,
      \tfrac{1}{2}\xi^{2\nu}\right) & \text{si}\ \boxed{n+1 = 2\nu} \\[6pt]
    \left(1, \xi, \ldots, \xi^{\nu}, [\xi^{\nu+1}, \lambda'_{\nu+1}],
      \tfrac{1}{2}\xi^{\nu+2}, \ldots, \tfrac{1}{2}\xi^{2\nu+1}\right) &
      \text{si}\ \boxed{n+1 = 2\nu+1}
  \end{array}
\]
avec $\xi\lambda'_{\nu+1} = 0$, $\lambda'^{\,2}_{\nu+1} = 0$.

\page{13}

\note{feuillet de tapuscrit du même exposé : la fin de la démonstration de
1.7 (le centralisateur $C$ d'un tore maximal est représentable et lisse, d'où
la constance locale de $s \mapsto \rho_n(s) = \dim C_s$) ; la convention
« $G$ de rang réductif localement constant » pour les conditions de 1.7 b) ;
le Corollaire 1.9 (si $\rho_r(s) = \rho_n(s)$, $\rho_r$ et $\rho_n$ sont
constants au voisinage de $s$, et le rang unipotent y est nul), avec sa
démonstration par 1.7 a) et l'inégalité $\rho_r \leqslant \rho_n$ ; le
Corollaire 1.10 ($G$ de rang réductif localement constant, le foncteur
$\mathcal{T} : (\mathrm{Sch})^{\circ} \to (\mathrm{Ens})$ des tores maximaux est
représentable par un préschéma lisse, séparé et de type fini sur $S$), avec
la démonstration du « type fini » par $\mathcal{T} \simeq G/\mathrm{Norm}_G(T)$}

\begin{itemize}
  \item les numéros « 1.7 », « 1.9 », « 1.10 » ont partout leur premier
  chiffre surchargé à l'encre en « 1 » ; les lettres $\rho$, $\mathcal{T}$,
  $\in$, $\leqslant$ sont de sa main ;
  \item « \struck{Corollaire} La démonstration de 1.7. est achevée » ;
  \item « représentable et lisse sur $S$ par \add{Exp. XI} 5.3. » ;
  \item « soit $s \in S$ tel que \add{$\rho_u(s) = 0$ i.e.} $\rho_r(s) =
  \rho_n(s)$ » ;
  \item « Il reste à vérifier que $\mathcal{T}$ est \struck{\ill{}} de type fini
  sur $S$ » ; « Mais alors $\mathcal{T}$ \struck{est} \add{isomorphe à}
  $G/\mathrm{Norm}_G(T)$ (comparer \add{Exp XI} 5.3. bis et 5.5.). ».
\end{itemize}

\page{14}

\struck{\ill{}}\note{en haut de la page, deux lettres barrées reliées par une
flèche}

\[
  \left\{
  \begin{array}{ll}
    H^i(E) = 0 & \text{si}\ i \neq 0, 2\nu - 2, 4\nu - 1 \\[4pt]
    H^{2(\nu-1)}(E) \simeq
      \left\{ \begin{array}{ll} 0, & \ell \neq 2 \\ \mathbf{Z}/2\mathbf{Z} & \text{si}\ \ell = 2 \end{array} \right. \\[10pt]
    H^{4\nu-1}(E) \simeq \mu(-2\nu)
  \end{array}
  \right.
  \qquad \boxed{n + 1 = 2\nu} \ \ n = 2\nu - 1
\]
\note{« $2\nu - 2$ » est surchargé sur « $2(\nu-1)$ » ; à côté, un
« \struck{si $i \neq 2(\nu-1)$} » biffé et une surcharge illisible}

\[
  \left\{
  \begin{array}{l}
    H^i(E) = 0 \quad \text{si}\ i \neq (0), 2\nu - 1, 2\nu + 1, 4\nu + 1 \\
    H^0(E) = \mu(0) \\
    H^{2\nu-1}(E) \simeq \mu(-\nu+1) \\
    H^{2\nu+1}(E) \simeq \mu(-\nu-1) \\
    H^{\uncertain{4\nu+1}}(E) \simeq \mu(-2\nu-1)
  \end{array}
  \right.
  \qquad \boxed{n + 1 = 2\nu + 1} \ \ n = 2\nu
\]
\note{les exposants des trois dernières lignes sont surchargés ; le « $0$ »
de la liste des exceptions est entouré}

\page{15}

\note{feuillet de tapuscrit du même exposé, dans le paragraphe 2 : la fin
d'une démonstration de semi-continuité (le sous-schéma ouvert et fermé $G'$,
fini sur $S$, de $G$ étale séparé, est un groupe constant $A_S$ sur le spectre
d'un anneau de valuation discrète complet ; le monomorphisme canonique $A \to B$
vers la fibre géométrique générique donne $f(s) \leqslant f(t)$), la
continuité de $f$ en $s$ équivalente à celle de $s \mapsto \mathrm{card}\, f(s)$
et à la finitude de $W$ au voisinage de $s$ ; puis le Lemme 2.3 : $G$ affine
lisse sur $k$ algébriquement clos, $R \subset T$ deux sous-tores ; alors $W(R)$
est isomorphe à un quotient d'un sous-groupe de $W(T)$}

\begin{itemize}
  \item « il contient un sous-schéma à la fois ouvert et fermé $G'$,
  \add{fini sur $S$,} tel que $G'_s = G_s$ (EGA II 6.2.6.) » ;
  \item « on a donc un \add{monomorphisme} \struck{\ill{}} canonique
  $A \to B$ » ;
  \item « indiquée dans 2.1. », « Lemme 2.3. » : le « 2 » surchargé à l'encre ;
  \item « indépendant de \struck{\ill{}} \add{la} structure de groupe
  \add{sur $W$}, et a été signalé après \add{Exp XI} 5.10. ; d'ailleurs sa
  démonstration \struck{\ill{}} se fait aisément par les arguments précédents,
  \add{en} utilisant le critère valuatif de propreté EGA II 7.3.8. » ;
  \item « comme dans \add{Exp XI} 5.9. » ;
  \item le dernier alinéa (« Dans la démonstration qui suit, il sera commode
  d'interpréter $G, R, T, C(R), N(R), C(T), N(T)$ de façon purement
  ensembliste … ») est encadré et barré.
\end{itemize}

\page{16}

\emph{Cohomologie à supports compacts} \struck{\ill{}} de $C, E, \widehat{E}$ …

La cohomologie à support compact de $E, \widehat{E}$ se déduit simplement par
dualité de la \add{coh. usuelle} \struck{\ill{}} de $E, \widehat{E}$ lorsque $X$
est lisse … celles de $\widehat{C}$ est identique à la cohomologie usuelle (car
$\widehat{C}$ est propre …).

Reste à déterminer la coh. à support\supplied{s} compact\supplied{s} de $C$.

Or on a \struck{$H^{*}_c(C') \Longleftarrow H^{*}_c(\ldots)$}
\[
  \cdots \to H^i(\widehat{C}) \to H^i(X) \to H^{i+1}_c(C) \to H^{i+1}(\widehat{C})
  \to H^{i+1}(X)
\]
avec $H^i(\widehat{C}) \simeq H^{i-2}(X)(-1)$ si $i \neq 0$, et
$H^{i+1}(\widehat{C}) \simeq H^{i-1}(X)(-1)$.

Comparant avec la \uncertain{suite exacte} \ill{} $H^{*}_c(\widehat{C})$ \ill{}
$\widehat{E} \subset \widehat{C}$, \ill{} \ill{} celle \ill{} \ill{} \ill{}
\uncertain{la coh. de} $E$ \ill{} \ill{} \ill{} détermine \ill{} :

prévoir, si $i \neq 0$, $H^i(E) \simeq H^{i+1}_c(C) \simeq H^{i+1}_a(C)$,
\struck{\ill{}}

\struck{\ill{} \ill{} \ill{} $E \subset C$,}
\[
  \boxed{H^{*}_a(C) \simeq H^{*}_c(C) \quad ??}
\]
\struck{$H^i_a(C) = H^i_c(C) =$}

\page{17}

\note{feuillet de tapuscrit du même exposé : la fin d'une remarque (tout tore
est contenu dans un tore maximal quand $S$ est artinien, ou local avec $G$
« réductif ») ; la Définition 1.5 des rangs réductif, nilpotent et unipotent de
$G$ sur un corps $k$ comme ceux de $G_{\bar k}$, invariants par extension du
corps de base ; la Remarque 1.6, qui construit sur le spectre d'un anneau de
valuation discrète un schéma en groupes affine et lisse de fibre générale
$\mathbf{G}_m$ et de fibre spéciale $\mathbf{G}_a$, sans tore maximal, où le
rang réductif des fibres n'est pas continu ; l'énoncé initial du Théorème 1.7
(notations $\rho_r(s)$, $\rho_n(s)$), qui se poursuit page 3}

\begin{itemize}
  \item « (on verra dans l'exposé \struck{ultérieur} \add{XIV} \add{que la
  réciproque est vraie} \struck{que la réciproque est vraie} lorsque $S$ est
  \add{artinien} \struck{\ill{}}, \struck{\ill{}} \add{ou} lorsque $S$ est un
  schéma local et $G$ est “réductif” » ;
  \item les numéros « 1.5 », « 1.2 », « 1.6 », « 1.7 » ont leur premier chiffre
  surchargé en « 1 » ; « Définition 1.5. Soit $G$ un groupe algébrique sur un
  corps $k$, \struck{\ill{}} » ;
  \item « avec l'extension de la base, \struck{\ill{}} les \struck{rangs}
  notions de rang introduites » ;
  \item « ne contient aucun tore sauf le tore trivial \add{$T$} (réduit au
  sous-groupe unité) » ; les $\rho$ de l'énoncé de 1.7 sont de sa main.
\end{itemize}

\page{19}

\emph{Cohomologie des quadriques}

$Q^n$ quadrique non singulière de dim $n$ sur $k$ sép. clos.

$Q^{n-1}$ section à l'infini.

$U^n = AQ^n = Q^n - Q^{n-1}$ \quad \uncertain{pour} $A$ coefficient annulé par
un $m$\note{lecture de « annulé par un $m$ » peu sûre}

\[
  H^i_c(U^n, A) \simeq
  \left\{
  \begin{array}{ll}
    0 & \text{si}\ i \neq n, 2n \\
    A \otimes \nu'_n & \text{si}\ i = n \\
    A \otimes \mu_m(-n) & \text{si}\ i = 2n
  \end{array}
  \right.
  \quad (\text{si}\ n \neq 0)
\]
\[
  \nu'_n = \check{\nu}_n \otimes \mu_m(-n) =
  \left\{
  \begin{array}{ll}
    \mu_m(-\tfrac{n}{2}) & n\ \text{pair} \\
    \mu_m(-\tfrac{n+1}{2}) & n\ \text{impair}
  \end{array}
  \right.
\]
\[
  \nu'_n = \check{\nu}_n = \mu_m(-\tfrac{n}{2}) \quad n\ \text{pair}, \qquad
  \nu'_n = \nu_n \otimes \mu_m(1) \quad n\ \text{impair}
\]
\note{les deux dernières égalités sont ajoutées d'une autre encre ; les indices
et accents des $\nu$ sont peu lisibles}

\marginal{si $n$ impair, \ill{} \ill{}, on ne sait \uncertain{déterminer} qu'au
signe près …}\note{écrit en biais dans la marge gauche, à côté d'une double
flèche verticale qui relie $H^i_c$ et $H^i$}

\uncertain{dualement}
\[
  H^i(U^n, A) \simeq
  \left\{
  \begin{array}{ll}
    0 & \text{si}\ i \neq 0, n \\
    A \otimes \nu_n & \text{si}\ i = 0 \\
    A & \text{si}\ i = n .
  \end{array}
  \right.
\]
\note{deux flèches croisées échangent les deux dernières lignes : il faut lire
$A$ pour $i = 0$ et $A \otimes \nu_n$ pour $i = n$ ; un « si $n \neq 0$ »
entouré les accompagne}
\[
  \nu_n \simeq \mu_m(-\tfrac{n}{2}) \ \text{si}\ n\ \text{pair}, \qquad
  \nu_n \simeq \mu_m(-\tfrac{n \pm 1}{2}) \ \text{si}\ n\ \text{impair}
\]
\note{le signe dans $n \pm 1$ est surchargé}

$H^i_c(U^n, A) \xrightarrow{\varphi} H^i(U^n, A)$ est \emph{nul} \emph{sauf} si
$n$ pair, $i = n$, (\uncertain{dans exc.~2}) \ill{} \ill{} \ill{} \ill{}
\ill{} \ill{} de $H^i_c(U^n, A)$ tel \uncertain{que}
$\langle \ell, \varphi(\ell) \rangle = 2$\note{une ligne entière biffée précède
ces trois derniers mots}

\[
  H^{i-1}(U^n) \to H^{i-2}(Q^{n-1})(-1) \to H^i(Q^n) \to H^i(U^n) \to
  H^{i-1}(Q^{n-1})(-1) \to \cdots
\]
\[
  H^{i-2}(Q^{n-1})(-1) \xrightarrow{\alpha_*} H^i(Q^n) \ \text{est}\
  \left|
  \begin{array}{ll}
    \text{surjectif} & \text{si}\ i \neq 0, n \\
    \text{injectif} & \text{si}\ i \neq 1, n+1
  \end{array}
  \right|
  \ \text{bijectif si}\ i \neq 0, \struck{\ill{}}\, n, n+1
\]

Reste à préciser ce qui se passe pour $i = 0, 1, n, n+1$ ??

Pour $i = 0, 1$ \quad $\alpha_*$ est \emph{nul} [$i = 1$, \struck{bij} car
bijectif, car les deux membres nuls) (si $\boxed{n \neq 1}$ \ill{} \ill{}
$\neq 0$]

Pour $i = n, n+1$ on a la \struck{\ill{}} suite exacte
\[
  0 \to H^{n-2}(Q^{n-1})(-1) \to H^n(Q^n) \to H^n(U^n) \to
  H^{n-1}(Q^{n-1})(-1) \to H^{n+1}(Q^n) \to 0
\]
\note{le « $(-1)$ » du premier terme semble barré, ou traversé par la flèche}

1) \emph{$n$ pair} $\geqslant 0$, la suite exacte \ill{} \ill{} \uncertain{se
réduit à} :
\[
  0 \to H^{n-2}(Q^{n-1})(-1) \to H^n(Q^n) \to H^n(U^n) \to 0
\]
donc $H^n(Q^n)$ est extension de $H^n(U^n)$ par $H^{n-2}(Q^{n-1})(-1)$ …

\page{20}

\note{feuillet de tapuscrit du même exposé, dans le paragraphe 2 : la
réduction au cas où $R_s$ est un tore maximal de $G_s$, l'égalité
$w(s) = w'(s)$ et les inégalités $w(t) \geqslant w'(t)$, $w'(t) \geqslant w'(s)$
pour $t$ voisin de $s$ ; le Lemme 2.2 ($W$ étale, séparé et quasi-fini sur $S$ ;
la fonction $f$ qui associe à $s$ la classe de la fibre géométrique de $W$ en
$s$, dans l'ensemble ordonné $E$ des classes de groupes finis, est
semi-continue inférieurement ; elle est constante au voisinage de $s$ si et
seulement si $s \mapsto \mathrm{card}\, f(s)$ l'est, si et seulement si $W$ est
fini au-dessus d'un voisinage de $s$) ; le début de sa démonstration
(réduction au cas affine noethérien, constructibilité, EGA IV 9, EGA
$0_{\mathrm{III}}$ 9.3.4, EGA II 7.1.7)}

\begin{itemize}
  \item « \struck{La deuxième} \add{Les} deux inégalités \struck{\ill{}} sont
  contenu\supplied{es} dans les deux lemmes suivants » ;
  \item « Lemme 2.2. », « introduit dans 2.1. » : le « 2 » surchargé ;
  « étale, \struck{quasi-fini et} séparé et quasi-fini sur $S$ » ;
  \item « qui est un \add{élément} \struck{ensemble} de l'ensemble ordonné $E$ » ;
  \item « semi-continue \supplied{inf}érieurement. Pour qu'elle soit constante au
  voisinage de $s \in S$, \struck{\ill{}} il faut et il suffit » ;
  \item « celui invoqué dans la démonstration de \add{Exp XI} 5.10., nous nous
  bornons à une esquisse de \add{la} démonstration (\add{qui est} du type le
  plus standard) » ;
  \item « et sor\supplied{it}es de EGA IV 9), et en vertu de EGA $0_{\mathrm{III}}$
  9.3.4. on est ramené \add{pour la semi-continuité} à prouver que si $t$ est
  une générisation de $s$, on a $f(t) \geqslant f(s)$ » ; « \struck{\ill{}}
  grâce à EGA II 7.1.7. ».
\end{itemize}

\end{document}
